% from ega4-1-fr.tex
\oldpage[IV-1]{235}

\begin{proposition}[1.6.2]
\label{IV.1.6.2-fr}
\begin{enumerate}
  \item[(i)] Toute immersion quasi-compacte et localement de présentation
  finie (en particulier toute immersion ouverte quasi-compacte) est de
  présentation finie.
  \item[(ii)] Le composé de deux morphismes de présentation finie est de
  présentation finie.
  \item[(iii)] Soient $X$, $Y$ deux $S$-préschémas, $f:X\to Y$ un
  $S$-morphisme de présentation finie. Alors, pour tout morphisme $S'\to S$,
  le morphisme $f_{(S')}:X_{(S')}\to Y_{(S')}$ est de présentation finie.
  \item[(iv)] Soient $f:X\to Y$, $f':X'\to Y'$ deux $S$-morphismes de
  présentation finie~; alors
  $f\times_Sf':X\times_SX'\to Y\times_SY'$ est de présentation finie.
  \item[(v)] Si le composé $g\circ f$ de deux morphismes est de présentation
  finie et si $g$ est quasi-séparé et localement de présentation finie, $f$
  est de présentation finie.
\end{enumerate}
\end{proposition}

Cela résulte immédiatement de (1.1.2), (1.2.2), (1.2.4) et (1.4.3).

Il résulte en particulier de (1.6.2, (iii)) que si $f$ est un morphisme de
présentation finie et $U$ un ouvert de $Y$, la restriction
$f^{-1}(U)\to U$ de $f$ est un morphisme de présentation finie. Inversement,
soit $(U_\alpha)$ un recouvrement de $Y$ par des ouverts affines, et supposons
que les restrictions $f^{-1}(U_\alpha)\to U_\alpha$ de $f$ soient des
morphismes de présentation finie~; il en résulte alors que $f$ est de
présentation finie, car $f$ est évidemment localement de présentation finie,
quasi-compact en vertu de (1.1.1) et quasi-séparé en vertu de (1.2.6). On peut
donc dire que la propriété pour un morphisme $f:X\to Y$ d'être de présentation
finie est \emph{locale sur $Y$}.

Si $X$ est un préschéma quasi-séparé, tout morphisme $f:X\to Y$ est
quasi-séparé (1.2.3). Donc, si $f$ est quasi-compact et localement de
présentation finie, $f$ est alors de présentation finie.

En particulier~:

\begin{corollary}[1.6.3]
\label{IV.1.6.3-fr}
Soit $\varphi:A\to B$ un homomorphisme d'anneaux. Pour que le morphisme
correspondant $f:\operatorname{Spec}(B)\to\operatorname{Spec}(A)$ soit de présentation finie, il faut et
il suffit que $\varphi$ fasse de $B$ une $A$-algèbre de présentation finie.
\end{corollary}

La nécessité de la condition résulte en effet de (1.4.6).

\begin{remark}[1.6.4]
\label{IV.1.6.4-fr}
Dans la définition (1.6.1), la condition (iii) n'est pas conséquence des deux
autres. Soit par exemple $Y$ un schéma affine dont l'espace sous-jacent ne
soit pas noethérien, et soit $U$ un ouvert non quasi-compact dans $Y$. Soit
$X$ le préschéma obtenu en recollant deux préschémas $Y_1$, $Y_2$ isomorphes à
$Y$ le long des ouverts $U_1$, $U_2$ correspondant à $U$ (0$_{\mathrm I}$,
4.1.7), de sorte que $X$ est réunion de deux ouverts isomorphes respectivement
à $Y_1$ et $Y_2$ (et que nous identifions à ces derniers), avec
$Y_1\cap Y_2=U$. En outre, il y a un morphisme $f:X\to Y$ qui coïncide dans
$Y_i$ avec l'isomorphisme donné $Y_i\to Y$ ($i=1,2$). Il est clair que $f$
vérifie la condition (i) de (1.6.1), étant un isomorphisme local~; il vérifie
aussi (ii), l'image réciproque d'un ouvert quasi-compact de $Y$ étant la
réunion de ses images dans $Y_1$ et $Y_2$~; mais comme $Y_1\cap Y_2$ n'est pas
quasi-compact, $f$ n'est pas quasi-séparé (1.2.7).
\end{remark}

\begin{proposition}[1.6.5]
\label{IV.1.6.5-fr}
Soient $X'$, $X''$ deux préschémas, $X=X'\amalg X''$ leur somme,
$f':X'\to Y$, $f'':X''\to Y$ deux morphismes, $f:X\to Y$ le morphisme qui
coïncide avec $f'$ dans $X'$ et avec $f''$ dans $X''$. Pour que $f$ soit de
présentation finie, il faut et il suffit que $f'$ et $f''$ le soient.
\end{proposition}

\par\medskip\noindent
\begin{proof}
Il suffit de montrer que pour que $f$ possède une des trois propriétés de la
définition

